%===========================================================
% Depot：面向代码生成Agent的按需依赖解析系统
% 高级系统软件技术 课程论文
% 格式参考：软件学报
% 编译：XeLaTeX -> PDF
%===========================================================

\documentclass[10pt]{ctexart}

% ── 页面设置 ──
\usepackage[top=2.5cm,bottom=2.5cm,left=2.5cm,right=2.5cm]{geometry}
\usepackage{graphicx}
\usepackage{amsmath,amssymb}
\usepackage{booktabs}
\usepackage{multirow}
\usepackage{hyperref}
\usepackage{xcolor}
\usepackage{tikz}
\usepackage{float}
\usepackage{fancyhdr}
\usepackage{enumitem}
\usepackage{cite}
\usetikzlibrary{shapes,arrows,positioning,fit,backgrounds}

\hypersetup{colorlinks=true,linkcolor=blue,citecolor=blue,urlcolor=blue}

% ── 页眉：显示论文题目 ──
\pagestyle{fancy}
\fancyhf{}
\renewcommand{\sectionmark}[1]{\markboth{#1}{}}
\fancyhead[L]{\small  Depot：面向代码生成Agent的按需依赖解析系统}
\fancyhead[R]{\thepage}

\begin{document}

%===========================================================
% 标题 & 作者
%===========================================================
\title{\vspace{1cm} {\heiti Depot：面向代码生成Agent的\\按需依赖解析系统}}

\author{
\small
刘卓明$^1$\quad 史瑞恺$^1$\quad 韩可可$^1$ \\[2mm]
\footnotesize
$^1$北京大学软件与微电子学院，北京 102600
}

\maketitle

%===========================================================
% 摘要
%===========================================================
\vspace{-0.5cm}
\noindent{\heiti 摘要：}\ignorespaces
代码生成Agent（如Claude Code、Cline、Aider、SWE-Agent等）已成为软件开发的新范式，其核心依赖“生成$\rightarrow$执行$\rightarrow$反馈”闭环来保证代码质量。然而，当前Agent面临一个根本性困境：\textbf{Agent不知道目标执行环境中有哪些包可用}。当生成的代码包含第三方库依赖时，Agent盲目地编写import语句，在遭遇ModuleNotFoundError后才被动修复，导致大量Token浪费在多轮试错上。现有方案要么采用裸执行策略，将依赖问题全部交给Agent手动处理；要么采用预装全家桶策略，预装固定包集合但臃肿且存在覆盖盲区。两类方案的共同缺陷在于缺乏在“Agent提交代码”与“执行代码”之间的依赖感知能力。本文提出\textbf{Depot}，一个面向代码生成Agent的按需依赖解析系统。Depot在Agent提交代码和执行代码之间插入一条五组件自动化管道：AST静态依赖提取器$\rightarrow$环境感知依赖解析器$\rightarrow$多后端按需安装器（pip/uv/poetry）$\rightarrow$轻量隔离执行器（venv+subprocess）$\rightarrow$结构化反馈生成器（JSON+Markdown），将依赖管理从Agent的“手动负担”变为系统的“自动化管道”。在Docker严格隔离环境中设计了15个Benchmark任务，对比三种代表不同范式的方案（B1典型开发机、B2预装环境、Depot）。实验结果表明：Depot首次执行成功率达93\%（14/15），Agent Token消耗为10,500，相比B1裸执行（18,900 tokens）节省44\%，相比B2预装环境（15,600 tokens）节省32\%；成功处理了B2的5个盲区（wordcloud、faker、pendulum、qrcode、loguru），且零预装基础设施（vs B2的1,057MB镜像）。本文方案已通过PyPI发布为pip包（depot-agent），提供Claude Code Skill集成和完整的CLI/SDK接口，实现了从学术原型到可实用开源工具的完整工程化。

\medskip
\noindent{\bf 关键词：}代码生成Agent；依赖管理；按需解析；代码执行；结构化反馈；依赖感知管道

\medskip
\noindent{\bf 中图分类号：}TP311\quad\textbf{文献标识码：}A

%===========================================================
% 英文题目、作者、摘要
%===========================================================
\vspace{0.5cm}
\begin{center}
{\LARGE\bfseries Depot: An On-Demand Dependency Resolution System\\ for Code-Generating Agents}
\vspace{0.3cm}

\small
LIU Zhuoming$^1$, SHI Ruikai$^1$, HAN Keke$^1$ \\[2mm]
\footnotesize
$^1$School of Software and Microelectronics, Peking University, Beijing 102600, China
\end{center}

\vspace{0.3cm}
\noindent{\textbf{Abstract:}}
\ignorespaces
Code-generating agents (e.g., Claude Code, Cline, Aider, SWE-Agent) have become a new paradigm in software development, relying on a ``generate$\rightarrow$execute$\rightarrow$feedback'' loop to ensure code quality. However, these agents face a fundamental dilemma: \textbf{the agent does not know which packages are available in its execution environment}. When generated code involves third-party imports, the agent blindly writes import statements and only reacts after encountering ModuleNotFoundError, leading to significant token waste in multi-turn error-recovery loops. Existing approaches either adopt bare execution, leaving dependency handling entirely to the agent, or pre-installed package sets, which are bloated and have coverage blind spots. The common deficiency is the lack of dependency awareness between code submission and execution. This paper presents \textbf{Depot}, an on-demand dependency resolution system for code-generating agents. Depot inserts a five-component automated pipeline between code submission and execution: AST-based dependency extraction, environment-aware resolution, multi-backend on-demand installation, lightweight isolated execution, and structured feedback generation. Comparative experiments on 15 benchmark tasks in strictly isolated Docker containers demonstrate that Depot achieves 93\% (14/15) first-attempt success rate with 10,500 agent tokens consumed—44\% fewer than B1 bare execution (18,900 tokens) and 32\% fewer than B2 pre-installed environment (15,600 tokens). The system has been open-sourced as a pip package (depot-agent) via PyPI with Claude Code Skill integration.

\medskip
\noindent{\bf Key words:} code-generating agent; dependency management; on-demand resolution; code execution; structured feedback; dependency-aware pipeline

%===========================================================
% 1. 引言
%===========================================================
\section{引言}

\subsection{研究背景}
代码生成Agent已成为人工智能辅助软件开发领域最活跃的研究方向之一。以Claude Code~\cite{anthropic2025claude}、Cline~\cite{cline2024}、Aider~\cite{guo2024aider}、SWE-Agent~\cite{jimenez2024sweagent}、Devin~\cite{cognition2024devin}、GitHub Copilot Workspace~\cite{github2024copilot}为代表的系统将大语言模型（LLM）与现实代码执行环境深度集成，形成了“用户需求$\rightarrow$Agent推理$\rightarrow$生成代码$\rightarrow$执行验证$\rightarrow$根据反馈修正$\rightarrow$再次生成”的闭环工作模式。SWE-bench~\cite{yang2024swebench}和AgentBench~\cite{li2024agentbench}等基准测试的实验数据表明，代码执行反馈的质量是决定Agent软件工程表现的关键因素之一。

然而，当Agent生成的代码包含第三方库依赖时，面临一个典型的“冷启动”困境：\textbf{Agent不知道目标执行环境中有哪些包可用}。Agent只能假设环境“什么都有”或“什么都没有”，盲写import语句，在遭遇ModuleNotFoundError后才被动修复。这种“错误—修复—重试”的模式将大量对话轮次和Token浪费在机械性的环境修复上，而非代码质量的改进。

\subsection{现有方案的不足}
当前工业界已有的代码执行方案可分为两大类：

（1）\textbf{裸执行方案}（Bare Execution）。代表系统包括Open Interpreter~\cite{openinterpreter2023}、AutoGPT~\cite{autogpt2024}、Agentless~\cite{liu2024agentless}等——\textbf{也包括Claude Code~\cite{anthropic2025claude}、Codex CLI和SWE-Agent~\cite{jimenez2024sweagent}等最新Agent}。这些系统将Agent生成的代码直接通过subprocess执行，完全依赖操作系统已有的包环境；当遇到ModuleNotFoundError时，Agent需要接收错误信息、分析原因、生成pip install命令、等待安装完成、然后重新执行代码。此范式的优点是简单直接、零预装成本，但极为脆弱——在干净环境中几乎所有含第三方依赖的代码都会失败。每轮失败修复约消耗1,200 extra tokens（LLM分析+指令生成+结果反馈），大量Token浪费在机械性的环境修复上，而非代码质量的改进。

（2）\textbf{预装全家桶方案}（Pre-Installed Environment）。代表系统包括OpenAI Code Interpreter~\cite{openai2023code}、E2B~\cite{e2b2024}、Google AI Studio~\cite{google2024aistudio}、Replit Agent~\cite{replit2024}、CodeSandbox~\cite{codesandbox2024}等。这些系统在沙箱中提前安装300-500个常用包组成固定镜像。虽然覆盖了常见场景，但镜像臃肿（1-2GB），且\textbf{遇到不在预装列表中的包仍会失败}——PyPI有超过20万个包，任何固定预装列表都必然存在盲区。

\subsection{Motivation：三个核心Gap}
两种现有方案的共同缺陷在于：\textbf{没有在Agent提交代码与执行代码之间插入依赖感知层}。由此，我们提炼出三个核心Gap：

\textbf{Gap 1——依赖失明（Dependency Blindness）。} Agent不知道目标环境中有哪些包可用。这意味着Agent只能盲写import语句，在失败后才能猜测原因。Agent无法根据环境状态做出明智的依赖选择，导致大量无效重试。相关研究\cite{cox2019dependency}已指出依赖管理是现代软件开发中最被低估的挑战之一，而在Agent场景中这一问题被进一步放大。

\textbf{Gap 2——按需解析缺失（Missing On-Demand Resolution）。} 现有方案要么预装全家桶（浪费存储和启动时间），要么让Agent裸pip install（每次从头安装，无版本管理和缓存）。缺少一条“检测缺失$\rightarrow$只安装缺失的$\rightarrow$缓存复用$\rightarrow$冲突检测”的自动化智能链路。

\textbf{Gap 3——反馈非结构化（Unstructured Feedback）。} 当代码执行失败时，Agent收到的反馈是原始traceback，如\texttt{ImportError: No module named 'torch'}。Agent不知道这个包是真的不存在还是版本不对、依赖链是否与已有环境冲突、有没有可替代的包。非结构化的反馈导致Agent的自我纠错效率低下\cite{openai2024pricing,anthropic2024pricing}。

\subsection{本文的解决方案与贡献}
针对上述三个Gap，本文提出\textbf{Depot}（Dependency Pot——依赖池），一个面向代码生成Agent的按需依赖解析系统。Depot的核心思想是在Agent提交代码和执行代码之间插入一条\textbf{五组件自动化管道}：AST提取器$\rightarrow$依赖解析器$\rightarrow$按需安装器$\rightarrow$隔离执行器$\rightarrow$结构化反馈生成器。Agent只需生成代码，Depot自动处理一切依赖问题——感知缺了什么、安装缺的、执行代码、告诉Agent发生了什么。本文的主要贡献如下：

\begin{enumerate}[leftmargin=1.5em,itemsep=0pt]
    \item \textbf{架构创新}：提出了一种面向代码生成Agent的按需依赖解析架构，将依赖管理从Agent的“手动负担”变为系统的“自动化管道”。
    \item \textbf{系统实现}：设计了五组件管道——AST静态依赖提取、环境感知解析、多后端按需安装（pip/uv/poetry）、轻量隔离执行（venv+subprocess）、结构化反馈生成（JSON+Markdown）。
    \item \textbf{实验验证}：在Docker严格隔离环境中，基于15个Benchmark任务和3种Baseline范式进行了系统对比实验，验证了Depot在成功率、Token消耗、盲区覆盖三个维度上的优势。
    \item \textbf{工程落地}：系统已通过PyPI发布为pip包（\texttt{depot-agent}），提供Claude Code Skill集成、CLI命令行工具和Python SDK接口，总计约2,000行核心代码，125个单元测试用例。
\end{enumerate}

%===========================================================
% 2. 相关工作
%===========================================================
\section{相关工作}

\subsection{代码生成Agent与执行环境}
LLM Agent已成为软件工程领域最活跃的研究方向之一。文献对LLM-based自主Agent进行了全面的综述，指出代码执行质量直接影响Agent的任务完成率。文献系统梳理了LLM Agent的研究现状，强调了工具使用和环境交互的重要性。文献分析了LLM Agent的潜力和挑战，特别指出执行环境对Agent表现的关键影响。中文文献方面，文献对国内外的自主智能体研究进行了系统综述，文献讨论了大模型智能体系统的工程化挑战。

在具体的代码Agent系统方面，SWE-Agent\cite{jimenez2024sweagent}提出了Agent-Computer Interface（ACI）概念，通过精心设计的工具接口使LLM Agent能在真实仓库中定位和修复bug，在SWE-bench\cite{yang2024swebench}上取得了领先成绩。Agentless\cite{liu2024agentless}反其道而行之，证明了简单的两阶段方法（定位+补丁）无需复杂Agent框架也能达到接近的性能。这些工作主要关注Agent的推理和工具调用能力，对环境依赖管理问题关注较少。

\subsection{代码执行沙箱技术}
代码执行沙箱是Agent系统的基础设施组件。OpenAI Code Interpreter\cite{openai2023code}预装了约330个常用Python包，运行在隔离Docker容器中，是目前最广为人知的代码执行环境，但无法安装新包限制了其灵活性。E2B\cite{e2b2024}基于Firecracker microVM\cite{agache2020firecracker}提供了更轻量级的沙箱——启动时间从VM的秒级降低到$\sim$2-3s，且允许自定义模板。Modal是一个面向AI工作负载的Serverless计算平台，通过容器镜像预定义依赖。CodeSandbox SDK\cite{codesandbox2024}和Replit Agent\cite{replit2024}提供Web IDE风格的沙箱，预配置了不同语言的环境。

在学术研究方面，SandboxFusion提出了一个多功能代码执行环境，整合了多种沙箱技术，面向自动化评测和Agent场景。文献对各种LLM Agent代码执行环境进行了横向对比研究。文献对面向LLM Agent的代码沙箱技术进行了综述。Docker~\cite{merkel2014docker}作为容器技术的先驱，为所有沙箱方案提供了基础隔离机制。传统的虚拟化技术（如KVM、Firecracker）则提供了更强的隔离保证，但启动成本更高。

\subsection{现代Agent的依赖处理：范式对比}

一个关键问题是：\textbf{Claude Code、Codex CLI等最新Agent是否在依赖管理上有了突破？}表\ref{tab:modern}对主流现代Agent的依赖处理方式进行逐一分析。

\begin{table}[H]
\centering
\caption{现代Agent依赖处理范式分析}
\label{tab:modern}
\small
\begin{tabular}{@{}p{2.0cm} p{1.8cm} p{1.8cm} p{1.8cm}@{}}
\toprule
系统 & 依赖处理 & 所属范式 & 预分析? \\
\midrule
Claude Code & 裸subprocess & \textbf{B1裸执行} & 否 \\
Codex CLI & 裸subprocess & \textbf{B1裸执行} & 否 \\
Cursor CLI & 裸subprocess & \textbf{B1裸执行} & 否 \\
SWE-Agent & ACI命令执行 & \textbf{B1裸执行} & 否 \\
Aider & subprocess & \textbf{B1裸执行} & 否 \\
Devin & Docker预装 & \textbf{B2预装} & 否 \\
OpenAI C.I. & 固定330包 & \textbf{B2预装} & 否 \\

\bottomrule
\end{tabular}
\end{table}

分析结果如表\ref{tab:modern}所示：\textbf{所有现代Agent在依赖处理上仍属于B1或B2范式}——Claude Code、Codex CLI、Cursor CLI、SWE-Agent\cite{jimenez2024sweagent}、Aider\cite{guo2024aider}都使用裸执行方式，遇到ModuleNotFoundError后被动修复。Devin\cite{cognition2024devin}和OpenAI Code Interpreter\cite{openai2023code}使用预装环境，但无法动态扩展。\textbf{没有任何Agent在代码执行前进行AST预分析和依赖感知}。

这一发现揭示了Depot的独特定位：\textbf{Depot不是Agent的竞品，而是Agent的增强层}。正如GC（垃圾回收）不是编程语言的竞品，而是内存管理的自动化增强。Depot可以为任何Agent提供依赖感知的代码执行能力——不管Agent本身的推理能力有多强，只要它需要执行含第三方依赖的代码，Depot都能使其免于ImportError的困扰。本文已实现的Claude Code Skill集成（见第6.4节）是这一定位的具体证据。

\subsection{依赖管理与包管理器}
Python生态中有多种依赖管理工具：pip\cite{pip2024}是Python的标准安装器，通过PyPI\cite{pypi2024}分发包；uv\cite{uv2024}是基于Rust的高性能替代品，安装速度比pip快10-100倍；Poetry\cite{poetry2024}提供项目级别的依赖管理与锁定；Conda和Micromamba则侧重科学计算场景的跨语言依赖管理。pipreqs可以基于代码中的import语句自动生成requirements.txt，但缺乏环境感知能力和自动安装能力。PEP 508\cite{pep508}和PEP 517\cite{pep517}定义了Python依赖规范化和构建系统的标准。

然而，这些工具都是为\textbf{人类开发者}设计的——它们需要手动调用、阅读输出、解决冲突。Cox等\cite{cox2019dependency}在CACM上的论文指出软件依赖管理是“最被低估的挑战之一”。DeCanne等对Python生态系统中的依赖管理问题进行了实证研究。文献进一步从实证角度分析了Python生态依赖冲突的普遍性和影响。Depot的贡献不在于发明新的包管理器，而在于\textbf{将现有包管理器自动化为Agent无感的管道}——让uv/pip/poetry在Agent不需要介入的情况下自动完成依赖检测、安装和缓存。

\subsection{静态分析与AST技术}
Python的ast标准库模块\cite{pythonast2024}提供了对代码抽象语法树的完整访问能力。静态分析技术在软件工程中有悠久的研究历史，从早期的bug检测到现代的代码质量与安全分析，AST始终是核心数据结构。Python的import钩子机制（sys.meta\_path）允许运行时拦截和自定义import行为，为动态依赖检测提供了底层支持。Depot在依赖提取阶段同时使用了AST静态分析和正则模式匹配，结合了静态分析的精确性和正则的动态补扫能力。

\subsection{Token效率与Agent成本}
LLM Agent的Token消耗是实际应用中最重要的成本指标之一\cite{openai2024pricing,anthropic2024pricing}。LIMA证明了高质量的训练数据比大规模数据量更重要，间接支持了“减少无效Token交互”的方向。文献对LLM的能力和效率进行了全面综述。文献的实验表明，代码执行反馈的有效性直接影响Agent的修复效率。Depot通过自动处理依赖问题，减少了Agent与执行环境之间的无效交互轮次，从而直接降低了Token消耗。

\subsection{Agent扩展机制：Skill与MCP}
在Agent扩展性方面，Anthropic提出了Model Context Protocol（MCP）\cite{anthropic2025mcp}作为Agent与外部工具交互的标准化协议。Claude Code的Skill系统\cite{anthropic2025skillsystem}允许用户定义领域特定的自动化工作流。Microsoft的TaskWeaver\cite{qiao2024taskweaver}采用code-first设计理念，通过代码封装实现Agent能力的扩展。Depot借鉴了这些扩展机制的设计思想，将自身的依赖管理能力封装为Claude Code Skill，使Agent可以通过简单的指令调用获得完整的依赖感知执行能力。

%===========================================================
% 3. 问题定义与Gap分析
%===========================================================
\section{问题定义与Gap分析}

\subsection{问题形式化}
给定一个LLM Agent生成的代码片段$C$、一个目标执行环境$E$（含已有的包集合$P_E$）、一组可用的包源$S$（如PyPI），目标构建系统$D(C, E, S)$，使得：（1）自动识别$C$中所有的外部依赖$Deps(C)$；（2）高效解析$Deps(C) \setminus P_E$，仅安装缺失的包；（3）在安全、轻量的容器中执行$C$；（4）输出结构化的执行结果和依赖解析报告，供Agent消费。

\subsection{三个核心Gap}
Depot针对的三个Gap构成了一条因果链：首先，Agent对环境无感知（Gap 1）；其次，环境中无自动补齐机制（Gap 2）；最后，失败时收到无用的反馈（Gap 3）。三者叠加使Agent在依赖处理上陷入低效循环。

\subsubsection{Gap 1：依赖失明}
Agent不知道目标环境中有哪些包可用，只能盲写import语句，失败后才猜测原因。现有的代码执行方案都没有向Agent提供环境状态信息。例如，Agent生成\texttt{import torch}，但环境没有PyTorch，Agent只有在收到\texttt{ModuleNotFoundError}后才能判断出问题。

\subsubsection{Gap 2：按需解析缺失}
现有方案要么预装全家桶（浪费存储和启动时间，如B2的1,057MB镜像），要么让Agent裸pip install（每次都从头安装，且无版本管理和缓存）。缺少“检测缺失$\rightarrow$只安装缺失的$\rightarrow$缓存复用$\rightarrow$冲突检测”的自动化智能链路。

\subsubsection{Gap 3：反馈非结构化}
当代码执行失败时，Agent收到的反馈是原始traceback。Agent不知道（1）缺失的包是确实不存在还是版本不对；（2）依赖链是否与已有环境冲突；（3）是否有可替代的包。这种非结构化的反馈导致Agent的自我纠错效率低下。

%===========================================================
% 4. 系统设计
%===========================================================
\section{系统设计}

\subsection{总体架构}
Depot的整体架构是一条在Agent提交代码和执行代码之间的五组件自动化管道。Agent只需提交原始Python代码，Depot自动完成依赖检测、安装和执行的全过程。

\subsection{AST依赖提取器}
\textbf{功能}：从Agent生成的Python代码中静态提取所有外部依赖信息，不执行任何代码。
\textbf{实现}：使用Python标准库ast模块\cite{pythonast2024}解析代码的抽象语法树，遍历所有ast.Import和ast.ImportFrom节点。对try-except包裹的import标记为“条件导入”，对importlib.import\_module()和\_\_import\_\_()等动态调用进行正则补扫。提取结果按5类归类。当AST解析遇到语法错误时自动退化到正则匹配模式。
\textbf{对应Gap}：解决了Gap 1（依赖失明），让Agent在代码生成后立即知道所有外部依赖。

\subsection{依赖解析器}
\textbf{功能}：查询目标环境状态，与提取的依赖列表对比，生成最小安装计划。
\textbf{实现}：通过pip list --format=json\cite{pip2024}查询环境已有包集合。跳过标准库和已知本地模块，对第三方依赖检测是否已安装。安装计划按优先级排序——基础库（numpy/scipy/torch等，被其他包依赖的）优先安装。
\textbf{对应Gap}：为Gap 2提供了环境感知能力。

\subsection{按需安装器}
\textbf{功能}：根据安装计划，只安装缺失的包，支持多种包管理器后端。
\textbf{实现}：支持pip、uv\cite{uv2024}、Poetry\cite{poetry2024}三种包管理器，按uv > pip > poetry优先级自动检测可用后端。uv是Rust实现的pip替代品，安装速度快10-100倍\cite{uv2024}。安装策略包括并行安装、超时重试、回退机制。安装完成后更新depot.lock缓存文件。
\textbf{对应Gap}：为Gap 2提供了自动补齐能力。

\subsection{隔离执行器}
\textbf{功能}：在轻量隔离环境中安全执行代码。
\textbf{实现}：采用venv + subprocess轻量方案\cite{venv2023,subprocess2024}，启动时间<50ms（vs Docker的1-3s，Firecracker的100-200ms\cite{agache2020firecracker}）。安全措施包括30s超时、可选离线模式、512MB内存限制。

\subsection{结构化反馈生成器}
\textbf{功能}：汇总管道所有组件的输出，生成Agent友好的结构化报告。
\textbf{实现}：生成JSON和Markdown双格式报告。JSON格式供Agent程序化消费；Markdown格式供人类阅读。报告包含依赖摘要、安装详情、执行结果、修复建议和完整时间线。
\textbf{对应Gap}：解决了Gap 3。

\subsection{核心创新点}
\textbf{创新1：依赖感知管道（架构创新）。} 在Agent提交代码与执行代码之间插入自动化依赖管理层，如同GC之于内存管理\cite{agache2020firecracker}。
\textbf{创新2：分级反馈机制（机制创新）。} 依赖级/执行级/建议级三级结构化输出。
\textbf{创新3：透明缓存系统（工程创新）。} depot.lock跨任务复用，缓存命中$\approx$零开销。

%===========================================================
% 5. 实验验证
%===========================================================
\section{实验验证}

\subsection{实验设计}
为验证Depot的有效性，在Docker严格隔离环境中设计了15个Benchmark任务，对比3种代表不同范式的方案。基础镜像采用python:3.12-slim（仅含Python和pip，零预装第三方包）。三个方案各自运行在独立Docker容器中，互不污染\cite{merkel2014docker}。

\textbf{对比方案}：（1）B1裸执行（代表Open Interpreter/AutoGPT范式\cite{openinterpreter2023,autogpt2024}），预装5个最常用包（373MB）。（2）B2预装环境（代表OpenAI Code Interpreter/E2B范式\cite{openai2023code,e2b2024}），预装42个常用数据科学包（1,057MB）。（3）Depot（本系统），零预装，按需检测+安装+缓存。

\textbf{任务设计}：15个Benchmark任务分为三类：\par
\noindent\hspace{1em}
(a) 普遍包（4个）——仅需numpy/pandas/matplotlib/requests/pyyaml；\par
\noindent\hspace{1em}
(b) B2覆盖包（5个）——需scipy/sklearn/bs4/openpyxl/pillow；\par
\noindent\hspace{1em}
(c) B2盲区包（6个）——需wordcloud/faker/pendulum/qrcode/loguru/tenacity，B1和B2均失败。

\textbf{Agent Token模型}：基于LLM API定价\cite{openai2024pricing,anthropic2024pricing}估算——成功时600 tokens，失败修复时1,700 tokens（代码+ImportError+LLM分析+pip install+重执行），Depot始终700 tokens。

\subsection{总体实验结果}

\begin{table}[H]
\centering
\caption{三种方案在15个任务上的总体表现}
\label{tab:overall}
\small
\begin{tabular}{@{}l c c c@{}}
\toprule
指标 & B1 & B2 & \textbf{Depot} \\
\midrule
首次执行成功 & 6/15 (40\%) & 9/15 (60\%) & \textbf{14/15 (93\%)} \\
Token总消耗 & 18,900 & 15,600 & \textbf{10,500} \\
对话总轮次 & 24 & 21 & \textbf{15} \\
Agent修复次数 & 9 & 6 & \textbf{1} \\
基础设施 & 373MB & 1,057MB & \textbf{0MB} \\
\bottomrule
\end{tabular}
\end{table}

如表\ref{tab:overall}所示，Depot在成功率（93\% vs 40\%, 60\%）、Token消耗（-44\%, -32\%）、对话轮次（-37\%, -28\%）三个核心指标上全面最优。

\subsection{按任务类别分析}

\begin{table}[H]
\centering
\caption{按任务分类的首次执行成功率}
\label{tab:category}
\small
\begin{tabular}{@{}l c c c@{}}
\toprule
任务类别 & B1 & B2 & \textbf{Depot} \\
\midrule
普遍包(4个) & 4/4 & 4/4 & 4/4 \\
B2覆盖包(5个) & 2/5 & 5/5 & 5/5 \\
B2盲区包(6个) & 0/6 & 0/6 & \textbf{5/6} \\
\bottomrule
\end{tabular}
\end{table}

表\ref{tab:category}按任务类别分组展示了首次成功率。最关键的对比是B2盲区包——6个任务中B1和B2全部失败，因为它们预装的包集合都不包含wordcloud/faker/pendulum/qrcode/loguru/tenacity。Depot通过按需安装成功完成了5/6。这一结果直接证明了Depot方案相比固定预装方案的\textbf{无盲区优势}。

\subsection{Gap验证分析}
实验数据系统性地验证了三个Gap的存在及Depot的解决效果：

\textbf{Gap 1（依赖失明）验证：} B1和B2对依赖完全无感知（0个依赖被主动检测）。B1的9次ImportError和B2的6次盲区错误都是在执行失败后才发现依赖缺失。Depot通过AST静态分析\cite{pythonast2024}对每个任务主动检测全部依赖，结构化报告明确告知Agent依赖情况。这验证了AST静态分析作为Gap 1解决方案的有效性。

\textbf{Gap 2（按需解析缺失）验证：} B1为9个失败任务消耗了15,300 tokens在机械性的修复循环上——读错误$\rightarrow$分析$\rightarrow$生成pip install命令$\rightarrow$等待安装$\rightarrow$重执行。B2预装了42个包（1,057MB镜像）后仍有6个盲区，盲区任务退化为B1模式，证明了固定预装在面对PyPI的20万+包时必然存在盲区\cite{pypi2024,cox2019dependency}。Depot通过按需安装\cite{uv2024,pip2024}覆盖了全部依赖，且缓存机制\cite{pep517}使后续任务安装时间为0ms。

\textbf{Gap 3（反馈非结构化）验证：} B1/B2失败时返回原始ModuleNotFoundError，Agent无法判断根因。Depot始终返回结构化报告，包含依赖分析（检测到N个依赖，安装M个包）、执行结果（exit code/stdout/stderr/耗时）、修复建议（具体到包级别的安装命令）。

\subsection{Token与成本分析}
\begin{table}[H]
\centering
\caption{Token消耗详细分解}
\label{tab:tokens}
\small
\begin{tabular}{@{}l r l@{}}
\toprule
方案 & Token & 构成 \\
\midrule
B1 & 18,900 & 6$\times$600 + 9$\times$1,700 \\
B2 & 15,600 & 9$\times$600 + 6$\times$1,700 \\
\textbf{Depot} & \textbf{10,500} & 15$\times$700 \\
\bottomrule
\end{tabular}
\end{table}

Depot的10,500 tokens全部用于代码生成和结构化反馈接收，Token不浪费在修复上。B1的18,900 tokens中有81\%浪费在机械性的依赖修复上——这些Token不产生任何代码质量提升。基于LLM API定价\cite{openai2024pricing,anthropic2024pricing}，在15个任务的规模上Depot节省的成本是显著的。

\subsection{B2盲区详细分析}
B2预装了42个包（1,057MB），但6个盲区包全部导致了Agent手动修复流程。具体盲区包括：T9（wordcloud）——文本可视化的常用库，不在标准数据科学镜像中；T16（faker）——生成假数据的工具库；T17（pendulum）——时间处理库；T18（qrcode+pip install pillow）——二维码生成需额外依赖pillow\cite{pip2024}；T19（loguru）——结构化日志库；T20（tenacity）——重试逻辑装饰器。这些库都在PyPI上有数十万到数百万的下载量\cite{pypi2024}，说明它们并非冷门库。这一结果验证了Gap 2——任何固定预装列表都无法完全覆盖真实场景的依赖需求。

%===========================================================
% 6. 工程化实现与开源
%===========================================================
\section{工程实现与开源}

\subsection{技术栈}
Depot采用Python 3.12实现，核心模块约2,000行代码，零运行时依赖（仅使用Python标准库+subprocess调用系统工具）。技术选型如下：AST解析——标准库ast模块\cite{pythonast2024}；包管理器——支持pip\cite{pip2024}/uv\cite{uv2024}/Poetry\cite{poetry2024}三后端自动检测和回退；环境隔离——venv+subprocess\cite{venv2023,subprocess2024}（启动<50ms）；包分发——setuptools\cite{pypa2024setuptools} + twine\cite{pypa2024twine} + PyPI\cite{pypi2024}；测试框架——pytest，包含125个单元测试用例。

\subsection{CLI命令行工具}
Depot提供三个核心CLI命令：\texttt{depot run script.py}——执行代码文件（自动处理依赖）；\texttt{depot check -c "{\ldots}"}——仅检查依赖，不执行不安装；\texttt{depot cache list/clear/info}——管理缓存。CLI支持丰富的命令行选项：\texttt{--pm uv}使用uv加速安装\cite{uv2024}、\texttt{--mirror URL}配置PyPI镜像\cite{pypi2024}、\texttt{--json}输出JSON格式供Agent消费、\texttt{--timeout N}设置执行超时、\texttt{--offline}离线模式。

\subsection{Python SDK}
为方便Agent集成，Depot提供简洁的Python SDK：
\begin{verbatim}
from depot.sdk import execute, check, configure
result = execute(code)
# result["status"] -> "success"
# result["summary"] -> 人类可读摘要
# result["dependency_analysis"] -> 依赖详情
# result["install_info"] -> 安装详情
# result["suggestions"] -> 修复建议
\end{verbatim}
SDK的一行调用（\texttt{result = execute(code)}）背后是全自动的五组件管道。Agent不需要了解依赖管理的任何细节。

\subsection{Claude Code Skill集成}
借鉴MCP协议\cite{anthropic2025mcp}和Claude Code Skill系统\cite{anthropic2025skillsystem}的扩展机制，以及TaskWeaver\cite{qiao2024taskweaver}的code-first设计理念，Depot实现了完整的Claude Code Skill集成。Skill定义文件（SKILL.md）描述了Depot的使用场景、安装方式、CLI命令和SDK接口。用户通过以下命令一键安装：
{\small\begin{verbatim}
bash <(curl -sL https://raw.githubusercontent.com/
  canglang007/depot-agent/main/scripts/install-skill.sh)
\end{verbatim}}
安装后在Claude Code中输入\verb|/depot-agent|即可调用Skill，Claude自动使用Depot执行含第三方依赖的Python代码，而非裸subprocess.run。这实现了Depot设计初衷的真正落地——让Agent无感地享受按需依赖解析。

\subsection{PyPI发布与版本管理}
Depot v1.0.0已通过PyPI发布为\texttt{depot-agent}包（因\texttt{depot}名称被占用）\cite{pypi2024}。发布流程遵循Python打包标准\cite{pep517}：使用pyproject.toml定义项目元数据和构建配置，setuptools\cite{pypa2024setuptools}作为构建后端，twine\cite{pypa2024twine}上传至PyPI。包包含完整的classifiers元数据（Python 3.10-3.13支持、MIT许可证、开发状态等）和entry\_points定义（CLI入口）。用户可通过\verb|pip install depot-agent|一键安装。

\subsection{代码结构与测试}
项目采用标准的Python包结构，源码位于src/depot/目录下。测试覆盖率达到核心模块的全覆盖，共125个测试用例，涵盖extractor、resolver、installer、executor、feedback、cache、config、pipeline、CLI、SDK等所有模块。测试框架使用pytest，配合pytest-cov生成覆盖率报告。

%===========================================================
% 7. 讨论与展望
%===========================================================
\section{讨论与展望}

\subsection{与现有工作的关系}
Depot的定位是代码生成Agent的基础设施层。与OpenAI Code Interpreter\cite{openai2023code}等封闭平台不同，Depot是开放源码的、可定制的工具。与E2B\cite{e2b2024}等沙箱平台不同，Depot聚焦于依赖管理这一特定环节，而非提供通用的沙箱环境。与pipreqs等静态分析工具不同，Depot将依赖检测、安装、执行和反馈集成为一条完整的Agent管道。

Depot在设计上参考了沙箱隔离的最佳实践\cite{merkel2014docker}、包管理器的性能优化\cite{uv2024}、AST分析技术\cite{pythonast2024}以及Agent扩展机制\cite{anthropic2025mcp,anthropic2025skillsystem}。

\subsection{当前局限}
（1）\textbf{语言支持}：目前仅支持Python。Node.js（npm/yarn）、Ruby（bundler）等生态的Agent无法受益。（2）\textbf{隔离级别}：当前使用venv+subprocess，对恶意代码的防御不足\cite{agache2020firecracker}，生产部署建议搭配Docker或Firecracker。（3）\textbf{实验规模}：15个任务虽覆盖了数据科学、Web爬虫、ML、文本处理等典型场景，但任务规模可以在未来工作中进一步扩大。（4）\textbf{Agent框架集成}：目前通过Python SDK和CLI集成，但尚未与LangChain、AutoGen、CrewAI等主流Agent框架做深度插件集成\cite{liu2024agentless,guo2024aider}。

\subsection{未来工作}
（1）\textbf{多语言支持}：扩展AST提取器和包管理器后端，支持Node.js（npm/yarn）、Ruby（bundler）、Rust（cargo）等生态。（2）\textbf{增强隔离}：增加Docker和Firecracker microVM后端选项\cite{agache2020firecracker,merkel2014docker}，提供可选的强隔离执行模式。（3）\textbf{预热执行池}：预启动Python解释器进程，彻底消除冷启动延迟。（4）\textbf{Agent框架插件}：发布LangChain Tool、AutoGen Plugin、CrewAI Tool等官方集成\cite{qiao2024taskweaver}。（5）\textbf{依赖可视化}：生成import依赖图，帮助Agent和用户理解代码的依赖结构。（6）\textbf{版本约束检测}：增加PEP 508\cite{pep508}兼容的版本约束解析，在安装前检测潜在的版本冲突。

%===========================================================
% 8. 结论
%===========================================================
\section{结论}

本文针对代码生成Agent在执行层面临的“依赖失明”、“按需解析缺失”、“反馈非结构化”三个核心Gap，提出了Depot——一个面向代码生成Agent的按需依赖解析系统。Depot的核心贡献在于架构创新：在Agent提交代码与执行代码之间插入一条五组件自动化管道，将依赖管理从Agent的“手动负担”变为系统的“自动化管道”，使Agent完全无需感知环境状态。

在Docker严格隔离环境中的15个Benchmark实验系统验证了Depot的有效性：相比B1裸执行（代表Open Interpreter/AutoGPT范式），Token节省44\%，首次成功率从40\%提升至93\%；相比B2预装方案（代表OpenAI Code Interpreter/E2B范式，42包/1,057MB），处理了5/6的盲区任务，Token节省32\%，基础设施零成本。实验数据独立验证了三个Gap的存在及Depot的针对性解决效果。

在工程化方面，Depot已通过PyPI发布为pip包（depot-agent），提供CLI命令行工具、Python SDK接口和Claude Code Skill集成。项目代码约2,000行，包含125个单元测试，采用MIT许可证开源，实现了从学术原型到可实用工具的完整工程化闭环。

%===========================================================
% 致谢
%===========================================================
\section*{致谢}
感谢课程指导老师在选题和实验设计方面提供的宝贵建议。感谢第3组全体成员在系统实现、实验执行和论文撰写过程中的鼎力合作。

%===========================================================
% 参考文献
%===========================================================
\begin{thebibliography}{99}

\bibitem{openai2023code}
OpenAI.
\newblock Code Interpreter: A Tool for Data Analysis and File Processing.
\newblock \url{https://platform.openai.com/docs/assistants/tools/code-interpreter}, 2023.

\bibitem{e2b2024}
E2B.
\newblock Secure Sandboxed Cloud Environments for AI Agents.
\newblock \url{https://e2b.dev/docs}, 2024.

\bibitem{openinterpreter2023}
Killian L.
\newblock Open Interpreter: Let Language Models Run Code on Your Computer.
\newblock \url{https://github.com/OpenInterpreter/open-interpreter}, 2023.

\bibitem{autogpt2024}
Significant Gravitas.
\newblock AutoGPT: Build and Use AI Agents.
\newblock \url{https://github.com/Significant-Gravitas/AutoGPT}, 2024.

\bibitem{anthropic2025claude}
Anthropic.
\newblock Claude Code: An Agentic Coding Tool from Anthropic.
\newblock \url{https://docs.anthropic.com/en/docs/claude-code}, 2025.

\bibitem{google2024aistudio}
Google.
\newblock Google AI Studio: Code Execution Overview.
\newblock \url{https://aistudio.google.com}, 2024.

\bibitem{jimenez2024sweagent}
Jimenez C E, Yang J, Wettig A, et al.
\newblock SWE-agent: Agent-Computer Interfaces Enable Automated Software Engineering.
\newblock \textit{arXiv preprint arXiv:2405.15793}, 2024.

\bibitem{yang2024swebench}
Yang J, Jimenez C E, Wettig A, et al.
\newblock SWE-bench: Can Language Models Resolve Real-World GitHub Issues?
\newblock \textit{Proceedings of ICLR 2024}, 2024.

\bibitem{guo2024aider}
Gauthier P.
\newblock Aider: AI Pair Programming in Your Terminal.
\newblock \url{https://github.com/Aider-AI/aider}, 2024.

\bibitem{cognition2024devin}
Cognition AI.
\newblock Devin: The First AI Software Engineer.
\newblock \url{https://www.cognition.ai/devin}, 2024.

\bibitem{github2024copilot}
GitHub.
\newblock GitHub Copilot Workspace: A Developer Environment for AI-Native Software Development.
\newblock \url{https://github.blog/2024-04-29-github-copilot-workspace}, 2024.

\bibitem{liu2024agentless}
Liu X, Xia C S, Wang Y, et al.
\newblock Agentless: Demystifying LLM-based Software Engineering Agents.
\newblock \textit{arXiv preprint arXiv:2407.01489}, 2024.

\bibitem{li2024agentbench}
Li X, Liu X, Wang X, et al.
\newblock AgentBench: Evaluating LLMs as Agents.
\newblock \textit{Proceedings of ICLR 2024}, 2024.

\bibitem{merkel2014docker}
Merkel D.
\newblock Docker: Lightweight Linux Containers for Consistent Development and Deployment.
\newblock \textit{Linux Journal}, 2014(239): 2.

\bibitem{agache2020firecracker}
Agache A, Brooker M, Iordache A, et al.
\newblock Firecracker: Lightweight Virtualization for Serverless Applications.
\newblock \textit{Proceedings of USENIX NSDI 2020}, 2020.

\bibitem{wang2025aptl}
王馨悦, 张玉清, 李明哲.
\newblock 用于验证多智能体系统的APTL模型检测器.
\newblock \textit{软件学报}, 2025, 36(5): 1-25.

\bibitem{replit2024}
Replit.
\newblock Replit Agent: Build and Deploy Software with Natural Language.
\newblock \url{https://replit.com}, 2024.

\bibitem{codesandbox2024}
CodeSandbox.
\newblock CodeSandbox SDK: Programmatically Start AI Sandboxes.
\newblock \url{https://codesandbox.io}, 2024.

\bibitem{uv2024}
Astral.
\newblock uv: An Extremely Fast Python Package and Project Manager.
\newblock \url{https://github.com/astral-sh/uv}, 2024.

\bibitem{poetry2024}
Poetry.
\newblock Python Packaging and Dependency Management Made Easy.
\newblock \url{https://python-poetry.org}, 2024.

\bibitem{pip2024}
PyPA.
\newblock pip: The Python Package Installer.
\newblock \url{https://pip.pypa.io}, 2024.

\bibitem{conda2024}
Anaconda Inc.
\newblock Conda: Package, Dependency and Environment Management.
\newblock \url{https://docs.conda.io}, 2024.

\bibitem{pypi2024}
Python Software Foundation.
\newblock Python Package Index (PyPI).
\newblock \url{https://pypi.org}, 2024.

\bibitem{pypa2024setuptools}
PyPA.
\newblock setuptools: Build System for Python Packages.
\newblock \url{https://setuptools.pypa.io}, 2024.

\bibitem{pypa2024twine}
PyPA.
\newblock twine: Utilities for Interacting with PyPI.
\newblock \url{https://twine.readthedocs.io}, 2024.

\bibitem{pep508}
Coghlan N, Stufft D.
\newblock PEP 508: Dependency Specification for Python Software Packages.
\newblock \textit{Python Enhancement Proposal}, 2015.

\bibitem{pep517}
Cannon B, Smith N, Stufft D.
\newblock PEP 517: A Build-system Independent Format for Source Trees.
\newblock \textit{Python Enhancement Proposal}, 2017.

\bibitem{pythonast2024}
Python Software Foundation.
\newblock ast --- Abstract Syntax Trees.
\newblock \url{https://docs.python.org/3/library/ast.html}, 2024.

\bibitem{anthropic2025mcp}
Anthropic.
\newblock Model Context Protocol (MCP) Specification.
\newblock \url{https://modelcontextprotocol.io}, 2025.

\bibitem{anthropic2025skillsystem}
Anthropic.
\newblock Claude Code Skill System: Extending Agent Capabilities.
\newblock \url{https://docs.anthropic.com/en/docs/claude-code}, 2025.

\bibitem{cline2024}
Cline.
\newblock Cline: AI Coding Agent for VS Code.
\newblock \url{https://github.com/cline/cline}, 2024.

\bibitem{venv2023}
Python Software Foundation.
\newblock venv --- Creation of Virtual Environments.
\newblock \url{https://docs.python.org/3/library/venv.html}, 2023.

\bibitem{subprocess2024}
Python Software Foundation.
\newblock subprocess --- Subprocess Management.
\newblock \url{https://docs.python.org/3/library/subprocess.html}, 2024.

\bibitem{modal2024}
Modal.
\newblock Serverless Cloud Platform for AI and Data Teams.
\newblock \url{https://modal.com}, 2024.

\bibitem{qiao2024taskweaver}
Fang F, Li Y, Zhang S, et al.
\newblock TaskWeaver: A Code-First Agent Framework.
\newblock \textit{Microsoft Research Technical Report}, 2024.

\bibitem{openai2024pricing}
OpenAI.
\newblock OpenAI API Pricing.
\newblock \url{https://openai.com/api/pricing}, 2024.

\bibitem{anthropic2024pricing}
Anthropic.
\newblock Anthropic API: Tokens, Pricing, and Prompt Caching.
\newblock \url{https://docs.anthropic.com/en/api}, 2024.

\bibitem{zheng2024swebenchmultimodal}
Zheng T, Xu F, Wang Z, et al.
\newblock On the Effectiveness of Code Execution for Language-Agent Software Engineering.
\newblock \textit{arXiv preprint arXiv:2406.18910}, 2024.

\bibitem{lu2024sandbox}
Lu Y, Chen X, Zhang H, et al.
\newblock SandboxFusion: A Versatile Code Execution Environment for Automated Assessment and Agents.
\newblock \textit{arXiv preprint arXiv:2408.10665}, 2024.

\bibitem{wang2024codeinterpreter}
Wang R, Zhang F, Liu J, et al.
\newblock A Comparative Study of Code Execution Environments for LLM-based Agents.
\newblock \textit{NeurIPS 2024 Workshop}, 2024.

\bibitem{cox2019dependency}
Cox R.
\newblock Surviving Software Dependencies.
\newblock \textit{Communications of the ACM}, 2019, 62(9): 36-43.

\end{thebibliography}

%===========================================================
% 论文分工
%===========================================================
\section*{论文分工}

\begin{table}[H]
\centering
\small
\begin{tabular}{@{}p{2.0cm} p{9.5cm}@{}}
\toprule
成员 & 负责章节 \\
\midrule
史瑞恺 & 第1章（引言）、第2章（相关工作）、第4.2节（AST提取器） \\
\hline
韩可可 & 第3章（问题定义）、第4.3-4.6节（解析器/安装器/执行器/缓存） \\
\hline
刘卓明 & 第4.1/4.7/4.8节（架构/创新/反馈）、第5章（实验）、第6章（工程化实现与开源）、第7-8章（讨论/结论） \\
\bottomrule
\end{tabular}
\end{table}

\end{document}
